\subsection{LOGARITHMIC MAPPING}

Lemma 1. Suppose that \(p \neq 0\) . Let \(G\) be a Lie group, \(G_1\) an open subgroup of \(G\) which is isomorphic to a standard group and \(x \in G\) . The following conditions are equivalent:

\begin{enumerate}
\def\labelenumi{(\roman{enumi})}
\item
  there exists a power of \(x\) which belongs to \(\mathbf{G}_1\) ;
\item
  there exists a strictly increasing sequence \((n_{i})\) of integers such that \(x^{n_{i}}\) tends to \(e\) as \(i\) tends to \(+\infty\) .
  (ii) \(\Rightarrow\) (i): obvious.
  (i) \(\Rightarrow\) (ii): suppose that \(y = x^{m} \in G_{1}\) . By Proposition 9 (i) of no. 5, \(y^{p^n}\) tends to \(e\) as \(n\) tends to \(+\infty\) , in other words \(x^{mp^n}\) tends to \(e\) as \(n\) tends to \(+\infty\) .
\end{enumerate}

PROPOSITION 10. Suppose that \(p \neq 0\) . Let \(G\) be a finite-dimensional Lie group. Let \(G_f\) be the set of \(x \in G\) for which there exists a strictly increasing sequence \((n_i)\) of integers such that \(x^{n_i}\) tends to \(e\) as \(i\) tends to \(+\infty\) .

\begin{enumerate}
\def\labelenumi{(\roman{enumi})}
\item
  \(G_{f}\) is open in G.
\item
  There exists one and only one mapping \(\psi\) of \(G_{f}\) into L(G) with the following properties:
\end{enumerate}

\begin{enumerate}
\def\labelenumi{(\alph{enumi})}
\item
  \(\psi (x^n) = n\psi (x)\) for all \(x\in \mathbf{G}_f\) and all \(n\in \mathbf{Z}\) ;
\item
  there exists an open neighbourhood \(\mathbf{V}\) of \(e\) in \(\mathbf{G}\) such that \(\psi |\mathbf{V}\) is the inverse mapping of an injective exponential mapping.
\end{enumerate}

\begin{enumerate}
\def\labelenumi{(\roman{enumi})}
\setcounter{enumi}{2}
\tightlist
\item
  The mapping \(\psi\) is analytic.
\end{enumerate}

There exists an open subgroup of \(\mathbf{G}\) which is isomorphic to a standard group (no. 3, Theorem 4). Assertion (i) then follows from Lemma 1.

Let U be an open subgroup of L(G) and \(\phi\colon\mathrm{U}\to\phi(\mathrm{U})\) an exponential mapping of G with the properties of Proposition 3 of no. 2. U can be assumed to be small enough for \(\phi(\mathrm{U})\subset\mathrm{G}_{f}\) . Let \(x\in\mathrm{G}_{f}\) . There exists \(m\in\mathbf{Z}-\{0\}\) such that \(x^{m}\in\phi(\mathrm{U})\) . The element \(\frac{1}{m}\phi^{-1}(x^{m})\) does not depend on the choice of \(m\) . For let \(m'\in\mathbf{Z}\) be such that \(x^{m'}\in\phi(\mathrm{U})\) . Then \(x^{mm'}\in\phi(\mathrm{U})\) and

\[
m ^ {\prime} \phi^ {- 1} (x ^ {m}) = \phi^ {- 1} (x ^ {m m ^ {\prime}}) = m \phi^ {- 1} (x ^ {m ^ {\prime}}),
\]

whence our assertion. Let \(\psi(x) = \frac{1}{m}\phi^{-1}(x^m)\) . Then \(\psi|\phi(U) = \phi^{-1}\) . On the other hand, if \(n \in \mathbf{Z}\) , then

\[
\psi (x ^ {n}) = \frac {1}{m} \phi^ {- 1} (x ^ {n m}) = \frac {n}{m} \phi^ {- 1} (x ^ {m}) = n \psi (x).
\]

Hence \(\psi\) has properties (a) and (b) of the proposition. In a neighbourhood of \(x\) , \(\psi\) is composed of the mappings \(x \mapsto x^m\) , \(y \mapsto \phi^{-1}(y)\) and \(z \mapsto \frac{1}{m} z\) ; hence \(\psi\) is analytic on \(G_f\) .

Finally, let \(\psi'\) be a mapping of \(G_f\) into L(G) and V' a neighbourhood of \(e\) in \(G_f\) such that \(\psi'(x^n) = n\psi'(x)\) for \(x \in G_f\) and \(n \in \mathbf{Z}\) and such that \(\psi'|V'\) is the inverse mapping of an injective exponential mapping. Then \(\psi\) and \(\psi'\) coincide on a neighbourhood W of \(e\) . If \(x \in G_f\) , there exists \(n \in \mathbf{Z}\) such that \(x^n \in W\) . Then

\[
n \psi^ {\prime} (x) = \psi^ {\prime} (x ^ {n}) = \psi (x ^ {n}) = n \psi (x)
\]

and hence \(\psi = \psi'\) .

DEFINITION 2. The mapping \(\psi\) of Proposition 10 is called the logarithmic mapping of G and is denoted by \(\log_{\mathbb{G}}\) or simply log.

PROPOSITION 11. Suppose that \(p \neq 0\) . Let \(x, y\) be two permutable elements of \(G_f\) . Then \(xy \in G_f\) and \(\log(xy) = \log x + \log y\) .

The fact that \(xy \in G_f\) follows from Lemma 1. Let \(U\) be an open subgroup of the additive group \(L(G)\) and \(\phi: U \to \phi(U)\) an exponential mapping of \(G\) with the properties of Proposition 3 of no. 2; \(U\) can be assumed to be small enough for \(\log(\phi(U))\) to be the inverse mapping of \(\phi\) . For \(n \in \mathbf{Z} - \{0\}\) suitably chosen, \(x^n \in \phi(U)\) , \(y^n \in \phi(U)\) . Let \(u = \log x^n\) , \(v = \log y^n\) , whence \(x^n = \phi(u)\) , \(y^n = \phi(v)\) . By formula (2), \([u,v] = 0\) . The Hausdorff formula proves then that \(\phi(\lambda(u + v)) = \phi(\lambda u)\phi(\lambda v)\) for \(|\lambda|\) sufficiently small; hence, for every sufficiently large integer \(i\) ,

\[
\phi (p ^ {i} (u + v)) = \phi (p ^ {i} u) \phi (p ^ {i} v)
\]

that is

\[
p ^ {i} (\log x ^ {n} + \log y ^ {n}) = \log (x ^ {n p ^ {i}} y ^ {n p ^ {i}})
\]

or

\[
n p ^ {i} (\log x + \log y) = n p ^ {i} \log (x y).
\]

PROPOSITION 12. Suppose that \(p \neq 0\) . Let \(x \in G_f\) . The following conditions are equivalent:

\begin{enumerate}
\def\labelenumi{(\roman{enumi})}
\item
  \(\log x = 0;\)
\item
  x is of finite order in G.
\end{enumerate}

If there exists an integer \(n > 0\) such that \(x^n = e\) , it follows that

\[
n \log x = \log x ^ {n} = 0,
\]

whence \(\log x = 0\) . If \(\log x = 0\) , let \(\mathrm{V}\) be a neighbourhood of \(e\) in \(\mathbf{G}_f\) such that \(\log |\mathrm{V}\) is the inverse mapping of an injective exponential mapping. There exists an integer \(n > 0\) such that \(x^n \in \mathrm{V}\) ; the equality \(\log x^n = 0\) implies \(x^n = e\) .

PROPOSITION 13. Suppose that \(p \neq 0\) . If \(G\) is compact or standard, then \(G_f = G\) .

If G is standard, it suffices to use Lemma 1. Suppose that G is compact. Let \(x \in G\) and V be a neighbourhood of \(e\) in G. Let \(y\) be a limit point of the sequence \((x^n)_{n \geqslant 0}\) . For all \(n > 0\) , there exist two integers \(n_1, n_2\) such that \(n_1 \geqslant 2n_2 \geqslant n\) and \(x^{n_1} \in y\mathrm{V}, x^{n_2} \in y\mathrm{V}\) , whence \(x^{n_1 - n_2} \in V^{-1}\mathrm{V}\) and \(n_1 - n_2 \geqslant n\) . Hence \(x \in G_f\) .

COROLLARY. Suppose that K is locally compact. Then \(G_{f}\) is the union of the compact subgroups of G.

Let \(x \in G\) . If \(x\) belongs to a compact subgroup of \(G\) , then \(x \in G_f\) (Proposition 13). Suppose that \(x \in G_f\) . As \(K\) is locally compact, there exists an open subgroup \(G_1\) of \(G\) which is compact. Then there exists an integer \(m > 0\) such that \(x^m \in G_1\) . The closed subgroup \(G_2\) generated by \(x^m\) is contained in \(G_1\) and is therefore compact. Then \(x\) commutes with the elements of \(G_2\) and hence \(G_2 \cup xG_2 \cup \cdots \cup x^{m-1}G_2\) is a compact subgroup of \(G\) which contains \(x\) .

Example. Suppose that K is locally compact. Let U be the set of invertible elements of A; it is a compact open subgroup of the Lie group K*. Then \(U \subset (\mathbf{K}^{*})_{f}\) by Proposition 13; on the other hand, if \(x \in K^{*}\) is such that \(x \notin U\) , either \(x^{n}\) tends to 0 as n tends to \(+\infty\) , or \(x^{n}\) tends to 0 as n tends to \(-\infty\) ; hence \(U = (\mathbf{K}^{*})_{f}\) . The function \(\log_{K^{*}}\) is defined and analytic on \(U_{f}\) with values in \(L(\mathbf{K}^{*}) = \mathbf{K}\) , and such that \(\log_{\mathbf{K}^{*}}(xy) = \log_{\mathbf{K}^{*}}(x) + \log_{\mathbf{K}^{*}}(y)\) for all x, y in U; the elements x of U such that \(\log_{\mathbf{K}^{*}}(x) = 0\) are the roots of unity of K.

We again use the notation of nos. 3, 4 and 5.

PROPOSITION 14. Suppose that \(p \neq 0\) and that \(v\) is chosen such that \(v(p) = 1\) . Let \(G\) be a standard group and \(E(X)\) (resp. \(L(X)\) ) the expansion of the exponential function of \(G\) (resp. the logarithmic function of \(G\) ) as an integral series about 0.

\begin{enumerate}
\def\labelenumi{(\roman{enumi})}
\item
  The domain of absolute convergence (Differentiable and Analytic Manifolds, R, 4.1.3) of E contains the set \(\Delta\) of \(x \in G\) such that \(\omega(x) > \frac{1}{p - 1}\) . Let \(\mathbf{E}'\) denote the mapping defined on \(\Delta\) by this series. Then \(\mathbf{E}'\) is an exponential mapping of G and is an isomorphism of the manifold \(\Delta\) onto itself.
\item
  The domain of absolute convergence of L contains G. Let L' denote the mapping on G defined by this series. Then L' is the logarithmic mapping of G and the restriction of L' to \(\Delta\) is the inverse mapping of E'.
\item
  The mapping \(\mathbf{E}'\) is an isomorphism of \(\Delta\) , with the Hausdorff law, onto the subgroup \(\Delta\) of \(\mathbf{G}\) .
\end{enumerate}

Using the notation of § 5, nos. 3 and 4, \(E = \sum_{m \geqslant 1} \frac{\psi_{m,m}}{m!}\) (§ 5, no. 4, Proposition 3). As the coefficients \(c_{\alpha\beta\gamma}\) belong to A, \(\|\psi_{m,m}\| \leqslant 1\) (Differentiable and Analytic Manifolds, R, Appendix) \(K^{r}\) is assumed to have the norm

\[
\left\| \left(\lambda_ {1}, \dots , \lambda_ {r}\right) \right\| = \sup \left(\left| \lambda_ {1} \right|, \dots , \left| \lambda_ {r} \right|\right).
\]

By Chapter II, § 8, no. 1, Lemma 1, \(v(m!) \leqslant \frac{m - 1}{p - 1}\) . If \(\omega(x) > \frac{1}{p - 1}\) , we see that \(m\omega(x) - v(m!)\) tends to \(+\infty\) with \(m\) , whence

\[
\left\| \frac {\psi_ {m , m}}{m !} \right\| \| x \| ^ {m} \leqslant \frac {1}{| m ! |} \| x \| ^ {m}
\]

which tends to 0 as m tends to +∞ and

\[
\omega \left(\frac {\psi_ {m , m} (x)}{m !}\right) > \frac {m}{p - 1} - \frac {m - 1}{p - 1} = \frac {1}{p - 1} \quad \text { for } m \geqslant 1.
\]

Therefore \(\Delta\) is contained in the domain of absolute convergence of E and \(\mathrm{E}'(\Delta) \subset \Delta\) . Clearly \(\mathrm{E}'\) is an exponential mapping.

If \(L_{m}\) denotes the homogeneous component of \(L\) of degree \(m\) , Proposition 3 of §5, no. 4, proves that each coefficient of \(L_{m}\) is of the form

\[
a _ {1} + \frac {1}{2} a _ {2} + \dots + \frac {1}{m} a _ {m}
\]

with \(a_1, a_2, \ldots, a_m\) in A; but

\[
\inf \left(v (1), v \left(\frac {1}{2}\right), \dots , v \left(\frac {1}{m}\right)\right) = o (\log m) \quad \text { as   } m \text {   tends   to   } + \infty
\]

and

\[
\inf \left(v (1), v \left(\frac {1}{2}\right), \dots , v \left(\frac {1}{m}\right)\right) \geqslant v \left(\frac {1}{m !}\right) \geqslant - \frac {m - 1}{p - 1}.
\]

Therefore, if \(\omega(x) > 0\) , \(\|L_m\|. \|x\|^m\) tends to 0 as \(m\) tends to \(+\infty\) , so that G is contained in the domain of absolute convergence of L. On the other hand, if \(\omega(x) > \frac{m}{p-1}\) , then \(\omega(L_m(x)) > \frac{m}{p-1} - \frac{m-1}{p-1} = \frac{1}{p-1}\) for \(m \geqslant 1\) and hence \(L'(\Delta) \subset \Delta\) .

As the formal power series \(\mathrm{L}(\mathrm{E}(\mathrm{X}))\) and \(\mathrm{E}(\mathrm{L}(\mathrm{X}))\) are equal to \(\mathbf{X}\) , no. 4.1.5 of Differentiable and Analytic Manifolds, \(\mathbf{R}\) , proves that

\[
\mathrm{L} ^ {\prime} (\mathrm{E} ^ {\prime} (x)) = \mathrm{E} ^ {\prime} (\mathrm{L} ^ {\prime} (x)) = x
\]

for \(x \in \Delta\) . Hence \(E'\) is an isomorphism of the manifold \(\Delta\) onto itself and the inverse isomorphism is the restriction of \(L'\) to \(\Delta\) .

\(\mathrm{L}(\mathrm{X}^{[n]}) = n\tilde{\mathrm{L}}(\mathrm{X})\) for \(n\) an integer \(> 0\) (cf.~§ 5, no. 4). As \(G\) is contained in the domain of absolute convergence of \(L\) and \(X^{[n]}\) , therefore \(L'(x^n) = nL'(x)\) for all \(x \in G\) . The relation \(L'| \Delta = E'^{-1}\) implies that \(L'(x^n) = \log x^n\) for \(n\) sufficiently large. Hence \(L'(x) = \log x\) . We have thus proved (i) and (ii).

Let \(H = \sum_{r,s \geq 0} H_{r,s}\) be the Hausdorff formal power series and h the Hausdorff function relative to L(G). The domain of absolute convergence of \(\tilde{H}\) contains \(\Delta \times \Delta\) and \(h\) is defined on \(\Delta \times \Delta\) (Chapter II, § 8, Proposition 2). Then

\[
\mathrm{E} ^ {\prime} (x) \mathrm{E} ^ {\prime} (y) = \mathrm{E} ^ {\prime} (h (x, y))
\]

for \(x, y\) sufficiently close to 0 (§ 4, Theorem 4 (v)). Hence, in the notation of no. 3, Definition 1, the formal power series \(\mathbf{F}(\mathrm{E}(\mathbf{X}), \mathrm{E}(\mathbf{Y}))\) and \(\mathrm{E}(\mathrm{H}(\mathbf{X}, \mathbf{Y}))\) are equal. Let \(x, y\) be elements of \(\Delta\) . Then

\[
\begin{array}{l} \sup _ {m} \left\| \frac {\psi_ {m , m}}{m !} \right\| (\sup \| x \|, \| y \|) ^ {m} <   1 \\ \sup _ {r, s} \| H _ {r, s} \| \| x \| ^ {r} \| y \| ^ {s} <   | p | ^ {1 / (p - 1)} \end{array}
\]

by Chapter II, §8, formula (14). By Differentiable and Analytic Manifolds, R, 4.1.5, \(\mathrm{E}'(x)\mathrm{E}'(y)\) is obtained by substituting \(x\) for X and \(y\) for Y in

\[
\mathrm{F} (\mathrm{E} (\mathrm{X}), \mathrm{E} (\mathrm{Y}))
\]

and \(\mathrm{E}'(h(x,y))\) is obtained by substituting \(x\) for \(\mathbf{X}\) and \(y\) for \(\mathbf{Y}\) in \(\mathrm{E}(\mathrm{H}(\mathbf{X},\mathbf{Y}))\) . Hence \(\mathrm{E}'(x)\mathrm{E}'(y) = \mathrm{E}'(h(x,y))\) .

